
\section{另两个技术引理}
\label{sec:two-other-lemmas}

我们现在来证明\trGluing{}的唯一性之证明中所用的两条引理\footnote{同样重要，尽管颇为技术性。}。下面是第一条。

\emph{引理（\cref{thm:chi-n-w1p}）。}
\[
\lim_{n\to+\infty}\|\chi_{n}\|_{W^{1,p}}=0\text{。}
\]

\begin{proof}
我们有
\[
\int_{\RR\times S^{1}}\|\chi_{n}\|^{2}\,ds\,dt=\int_{\RR\times S^{1}}\langle\chi_{n},\chi_{n}\rangle\,ds\,dt=\int_{\RR\times S^{1}}\langle\zeta_{n},\chi_{n}\rangle\,ds\,dt\text{。}
\]
比如说，应用\cref{thm:w1p-product-inequality}，我们找到一个常数 $M>0$，使得
\[
\int_{\RR\times S^{1}}\langle\zeta_{n},\chi_{n}\rangle\,ds\,dt\le M\|\zeta_{n}\|_{L^{\infty}}\int_{\RR\times S^{1}}\|\chi_{n}\|\,ds\,dt\text{。}
\]
由于 $Y_{n}=i(\zeta_{n})$、$(\zeta_{n})_{j}=\langle Y_{n},Z_{j}\rangle$，并由\cref{thm:lift-exp-decay}，我们有 $\lim_{n\to+\infty}\|\zeta_{n}\|_{L^{\infty}}=0$。设 $\varepsilon>0$。从上面我们推出：存在一个 $N_{\varepsilon}$，使得对每个 $n\ge N_{\varepsilon}$，
\[
\int_{\RR\times S^{1}}\|\chi_{n}\|^{2}\,ds\,dt\le M\varepsilon\int_{\RR\times S^{1}}\|\chi_{n}\|\,ds\,dt\text{。}
\]
我们把 $\chi_{n}$ 写在 $W_{n}$ 的“\trCanonical”基中：
\[
\chi_{n}=\alpha_{n}\Bigl(\frac{\partial u}{\partial s}\,\#_{w_{n}}\,0\Bigr)+\beta_{n}\Bigl(0\,\#_{w_{n}}\,\frac{\partial v}{\partial s}\Bigr)\qquad\text{其中 }\alpha_{n},\beta_{n}\in\RR\text{。}
\]
由此推出
\[
\begin{aligned}
\int_{\RR\times S^{1}}\|\chi_{n}\|\,ds\,dt&=\int_{]-\infty,-1]\times S^{1}}\|\chi_{n}\|\,ds\,dt\\
&\quad+\int_{[-1,1]\times S^{1}}\|\chi_{n}\|\,ds\,dt+\int_{[1,+\infty[\times S^{1}}\|\chi_{n}\|\,ds\,dt\\
&=\int_{]-\infty,-1]\times S^{1}}|\alpha_{n}|\,\Bigl\|\frac{\partial u}{\partial s}(s+\rho_{n},t)\Bigr\|\,ds\,dt\\
&\quad+\int_{[-1,1]\times S^{1}}\|\chi_{n}\|\,ds\,dt+\int_{[1,+\infty[\times S^{1}}|\beta_{n}|\,\Bigl\|\frac{\partial v}{\partial s}(s-\nu(\rho_{n}),t)\Bigr\|\,ds\,dt\text{。}
\end{aligned}
\]
我们将要对这三项各给出一个上界。为了简化记号并缩短公式，我们令
\[
M_{n}=\sqrt{\int_{]-\infty,-1]\times S^{1}}|\alpha_{n}|^{2}\,\Bigl\|\frac{\partial u}{\partial s}(s+\rho_{n},t)\Bigr\|^{2}\,ds\,dt}\text{。}
\]
我们显然有
\[
M_{n}\le\|\chi_{n}\|_{L^{2}}\text{。}
\]
于是
\[
\begin{aligned}
&\int_{]-\infty,-1]\times S^{1}}|\alpha_{n}|\,\Bigl\|\frac{\partial u}{\partial s}(s+\rho_{n},t)\Bigr\|\,ds\,dt\\
&\qquad=M_{n}\frac{\displaystyle\int_{]-\infty,-1]\times S^{1}}\Bigl\|\frac{\partial u}{\partial s}(s+\rho_{n},t)\Bigr\|\,ds\,dt}{\displaystyle\sqrt{\int_{]-\infty,-1]\times S^{1}}\Bigl\|\frac{\partial u}{\partial s}(s+\rho_{n},t)\Bigr\|^{2}\,ds\,dt}}\\
&\qquad=M_{n}\frac{\displaystyle\int_{]-\infty,-1+\rho_{n}]\times S^{1}}\Bigl\|\frac{\partial u}{\partial s}(s,t)\Bigr\|\,ds\,dt}{\displaystyle\sqrt{\int_{]-\infty,-1+\rho_{n}]\times S^{1}}\Bigl\|\frac{\partial u}{\partial s}(s,t)\Bigr\|^{2}\,ds\,dt}}\text{。}
\end{aligned}
\]
上式最后一个分式的分子可以从上加以控制：
\[
\int_{]-\infty,-1+\rho_{n}]\times S^{1}}\Bigl\|\frac{\partial u}{\partial s}(s,t)\Bigr\|\,ds\,dt\le\Bigl\|\frac{\partial u}{\partial s}\Bigr\|_{L^{1}(\RR\times S^{1};\RR^{2n})}
\]
而其分母可以从下加以控制：
\[
\sqrt{\int_{]-\infty,-1+\rho_{n}]\times S^{1}}\Bigl\|\frac{\partial u}{\partial s}(s,t)\Bigr\|^{2}\,ds\,dt}\ge\Bigl\|\frac{\partial u}{\partial s}\Bigr\|_{L^{2}(]-\infty,0]\times S^{1};\RR^{2n})}\text{。}
\]
令
\[
K_{1}=\frac{\Bigl\|\dfrac{\partial u}{\partial s}\Bigr\|_{L^{1}(\RR\times S^{1};\RR^{2n})}}{\Bigl\|\dfrac{\partial u}{\partial s}\Bigr\|_{L^{2}(]-\infty,0]\times S^{1};\RR^{2n})}}\text{，}
\]
于是我们有
\[
\int_{]-\infty,-1]\times S^{1}}|\alpha_{n}|\,\Bigl\|\frac{\partial u}{\partial s}(s+\rho_{n},t)\Bigr\|\,ds\,dt\le K_{1}\|\chi_{n}\|_{L^{2}}\text{。}
\]
与 $K_{1}$ 一样，下面出现的所有常数 $K_{i}>0$ 都不依赖于 $n$。以同样的方式，对第三项我们有
\[
\int_{[1,+\infty[\times S^{1}}|\beta_{n}|\,\Bigl\|\frac{\partial v}{\partial s}(s-\nu(\rho_{n}),t)\Bigr\|\,ds\,dt\le K_{2}\|\chi_{n}\|_{L^{2}(\RR\times S^{1};\RR^{2n})}
\]
而中间那一项只需用\trCauchy--\trSchwarz{}不等式便可从上加以控制：
\[
\int_{[-1,1]\times S^{1}}\|\chi_{n}\|\,ds\,dt\le K_{3}\|\chi_{n}\|_{L^{2}([-1,1]\times S^{1};\RR^{2n})}\le K_{3}\|\chi_{n}\|_{L^{2}(\RR\times S^{1};\RR^{2n})}\text{。}
\]
我们于是得到
\[
\|\chi_{n}\|_{L^{1}}\le K_{4}\|\chi_{n}\|_{L^{2}}\text{，}
\]
从而对 $n\ge N_{\varepsilon}$，我们有
\[
\|\chi_{n}\|_{L^{2}}^{2}\le MK_{4}\varepsilon\,\|\chi_{n}\|_{L^{2}}\text{，}\qquad\text{从而 }\|\chi_{n}\|_{L^{2}}\le MK_{4}\varepsilon\text{，}
\]
因此
\[
\lim_{n\to+\infty}\|\chi_{n}\|_{L^{2}}=0\qquad\text{并且最终 }\lim_{n\to+\infty}\|\chi_{n}\|_{L^{1}}=0\text{。}
\]

回到 $\chi_{n}$ 的分解，由此推出\footnote{译注：下式积分区域原书印作 $]-\infty,1]\times S^{1}$；按被积函数是 $\chi_{n}$ 分解中来自 $u$ 的那一项（$u$ 型分量只在 $s\le-1$ 的区域上出现），且其后结论 $\lim_{n\to+\infty}\alpha_{n}=0$ 也只对该区域有意义，疑应为 $]-\infty,-1]\times S^{1}$。此处径改。——译者注}
\[
\lim_{n\to+\infty}\int_{]-\infty,-1]\times S^{1}}|\alpha_{n}|\,\Bigl\|\frac{\partial u}{\partial s}(s+\rho_{n},t)\Bigr\|=0\text{，}
\]
从而 $\lim_{n\to+\infty}\alpha_{n}=0$，同样地 $\lim_{n\to+\infty}\beta_{n}=0$。于是
\[
\|\chi_{n}\|_{W^{1,p}}\le|\alpha_{n}|\,\Bigl\|\frac{\partial u}{\partial s}\,\#_{w_{n}}\,0\Bigr\|_{W^{1,p}}+|\beta_{n}|\,\Bigl\|0\,\#_{w_{n}}\,\frac{\partial v}{\partial s}\Bigr\|_{W^{1,p}}
\]
趋于 $0$，因为我们所用的 $W_{n}$ 之基中那两个向量的 $W^{1,p}$ 范数显然是\trBounded{}的（关于 $n$ 一致\trBounded）。这就结束了\cref{thm:chi-n-w1p}的证明。
\end{proof}

我们要在这里证明的第二条引理（它用于\trGluing{}的唯一性的证明）是下面这条。

\emph{引理（\cref{thm:yn-upper-bound}）。}存在正常数 $K_{1}$ 与 $K_{2}$，使得对每个 $n\in\NN$，
\[
\|Y_{n}\|_{W^{1,p}}\le K_{1}\,\frac{\|\chi_{n}\|_{W^{1,p}}+\|Y_{n}\|_{L^{\infty}}+\|\widetilde{\scr{F}}(0)\|_{L^{p}}}{1-K_{2}\|Y_{n}\|_{L^{\infty}}}\text{。}
\]

我们当然使用在\cref{thm:yn-upper-bound}附近定义的那些记号。我们提醒一下，$\zeta_{n}\in W^{1,p}(\RR\times S^{1};\RR^{2n})$ 可以分解成 $\chi_{n}+\omega_{n}$，其中 $\chi_{n}$ 具有 $\alpha\,\#_{w_{n}}\,\beta$ 的形式（$\alpha\in\operatorname{Ker}L^{u}$，$\beta\in\operatorname{Ker}L^{v}$），而 $\omega_{n}$ 与这种形式的诸向量\trOrthogonal（参见\cref{sec:uniqueness-gluing}）。

\begin{proof}
下面出现的诸常数 $C_{i}$ 都不依赖于 $n$。我们有
\[
\begin{aligned}
\|Y_{n}\|_{W^{1,p}}=\|i(\zeta_{n})\|_{W^{1,p}}\le C_{1}\,\|\zeta_{n}\|_{W^{1,p}}=C_{1}\,\|\chi_{n}+\omega_{n}\|_{W^{1,p}}\\
\le C_{1}\left(\|\chi_{n}\|_{W^{1,p}}+\|\omega_{n}\|_{W^{1,p}}\right)\text{，}
\end{aligned}
\]
其中我们用到了 $\|i\|$ \trBounded（不依赖于 $n$）这一事实，这已在\cref{sec:operator-frho}中证明。我们把\cref{thm:lrho-lower-bound}应用于 $\omega_{n}\in W_{n}^{\perp}$，便得到
\[
\begin{aligned}
\|Y_{n}\|_{W^{1,p}}&\le C_{1}\left(\|\chi_{n}\|_{W^{1,p}}+\|\omega_{n}\|_{W^{1,p}}\right)\\
&\le C_{1}\|\chi_{n}\|_{W^{1,p}}+C_{2}\|L(\omega_{n})\|_{L^{p}}\\
&=C_{1}\|\chi_{n}\|_{W^{1,p}}+C_{2}\|L(\zeta_{n}-\chi_{n})\|_{L^{p}}\\
&\le C_{1}\|\chi_{n}\|_{W^{1,p}}+C_{2}\|L(\zeta_{n})\|_{L^{p}}+C_{2}\|L(\chi_{n})\|_{L^{p}}\text{。}
\end{aligned}
\]
\trLinearOperator{} $L$ 的范数从上被一个不依赖于 $n$ 的常数控制。事实上，$L=(dF)_{0}$，这里 $F=p\circ(\operatorname{Id}\times\widetilde{\scr{F}})\circ i$。对于出现在这个复合中的诸算子，我们已经验证了\cref{thm:composition-rho-derivative-estimate}的条件 (2) 与 (4)，而这正是我们断定该复合之\trDifferential{}的范数 $\|(dF)_{0}\|$ 一致\trBounded{}所需要的。然而我们必须强调，这些结果是对\trPregluing{} $w_{\rho}$ 得到的，而不是对形如 $w_{\nu,\rho}$ 的\trPregluing{}得到的。诸读者可以验证，在更一般的情形，论证在一切细节上都是类似的。

让我们回到刚才得到的关于 $Y_{n}$ 的范数的估计。从上面我们推出
\[
\begin{aligned}
\|Y_{n}\|_{W^{1,p}}&\le C_{1}\|\chi_{n}\|_{W^{1,p}}+C_{2}\|L(\zeta_{n})\|_{L^{p}}+C_{3}\|\chi_{n}\|_{W^{1,p}}\\
&=(C_{1}+C_{3})\|\chi_{n}\|_{W^{1,p}}+C_{2}\|L(\zeta_{n})\|_{L^{p}}\text{。}
\end{aligned}
\]
由\cref{thm:chi-n-w1p}，我们知道第一项趋于 $0$。我们尚须从上控制 $L(\zeta_{n})$ 的 $L^{p}$ 范数。我们有
\[
\begin{aligned}
L(\zeta_{n})=(dF)_{0}(\zeta_{n})&=\Bigl(d\bigl(p\circ(\operatorname{Id}\times\widetilde{\scr{F}})\circ i\bigr)\Bigr)_{0}(\zeta_{n})\\
&=\Bigl(d\bigl(p\circ(\operatorname{Id}\times\widetilde{\scr{F}})\bigr)\Bigr)_{0}i(\zeta_{n})\\
&=\Bigl(d\bigl(p\circ(\operatorname{Id}\times\widetilde{\scr{F}})\bigr)\Bigr)_{0}(Y_{n})\\
&=(D_{1}p)_{(0,\widetilde{\scr{F}}(0))}(Y_{n})+(D_{2}p)_{(0,\widetilde{\scr{F}}(0))}(d\widetilde{\scr{F}})_{0}(Y_{n})\text{。}
\end{aligned}
\]
映射 $p=p_{\rho}$（它现在依赖于 $n$），
\[
p:W^{1,p}(\RR\times S^{1};\RR^{m})\times L^{p}(\RR\times S^{1};\RR^{m})\longrightarrow L^{p}(\RR\times S^{1};\RR^{2n})\text{，}
\]
由公式
\[
p\bigl((X,Y)\bigr)_{i}(s,t)=P\bigl(w_{n}(s,t),X(s,t),Z_{i}(s,t),Y(s,t)\bigr)\text{，}
\]
来定义，其中映射
\[
P:W\times\RR^{m}\times\RR^{m}\times\RR^{m}\longrightarrow\RR
\]
\footnote{译注：此陪域原书印作 $\RR^{2n}$；按上一式 $p$ 的第 $i$ 个分量即 $P\bigl(w_{n}(s,t),X(s,t),Z_{i}(s,t),Y(s,t)\bigr)$ 为标量，此处 $P$ 正是\cref{sec:operator-frho}中所定义的那个逐分量使用的实值映射（彼处陪域即 $\RR$），疑应作 $\RR$。此处径改。——译者注}对最后一个变元是线性的。我们已经证明过（在\cref{thm:df-drho-estimate}的证明中）：对所有的 $v,w\in W^{1,p}(\RR\times S^{1};\RR^{m})$，我们有
\[
\bigl\|(D_{1}p)_{(0,\widetilde{\scr{F}}(0))}(v)\bigr\|_{L^{p}}\le C_{4}\|v\|_{L^{p}}\qquad\text{并且}\qquad\bigl\|(D_{2}p)_{(0,\widetilde{\scr{F}}(0))}(w)\bigr\|_{L^{p}}\le C_{5}\|w\|_{L^{p}}\text{。}
\]
注意到 $\widetilde{\scr{F}}(0)=\scr{F}(w_{n})$ 在 $s$ 落于 $[-1,1]$ 之外时为零（因为此时 $w_{n}$ 是一条\trFloer\trTrajectory），我们可以把第一个估计稍加改进。由于 $p$ 对它的第二个变元是线性的，由此推出：对每个 $v\in W^{1,p}(\RR\times S^{1};\RR^{m})$，我们有
\[
\forall\,(s,t)\in(\RR-[-1,1])\times S^{1}\text{，}\quad(D_{1}p)_{(0,\widetilde{\scr{F}}(0))}v(s,t)=0\text{。}
\]
\cref{thm:df-drho-estimate}的证明中所用的论证于是给出
\[
\bigl\|(D_{1}p)_{(0,\widetilde{\scr{F}}(0))}v\bigr\|_{L^{p}}\le C_{4}\|v\|_{L^{p}([-1,1]\times S^{1};\RR^{m})}\text{。}
\]
利用这些评注，我们给出 $\|L(\zeta_{n})\|_{L^{p}}$ 的一个控制：
\[
\begin{aligned}
\|L(\zeta_{n})\|_{L^{p}}&\le\bigl\|(D_{1}p)_{(0,\widetilde{\scr{F}}(0))}(Y_{n})\bigr\|_{L^{p}}+\bigl\|(D_{2}p)_{(0,\widetilde{\scr{F}}(0))}(d\widetilde{\scr{F}})_{0}(Y_{n})\bigr\|_{L^{p}}\\
&\le C_{4}\|Y_{n}\|_{L^{p}([-1,1]\times S^{1};\RR^{m})}+C_{5}\bigl\|(d\widetilde{\scr{F}})_{0}(Y_{n})\bigr\|_{L^{p}(\RR\times S^{1};\RR^{m})}\\
&\le C_{6}\|Y_{n}\|_{L^{\infty}([-1,1]\times S^{1};\RR^{m})}+C_{5}\bigl\|(d\widetilde{\scr{F}})_{0}(Y_{n})\bigr\|_{L^{p}(\RR\times S^{1};\RR^{m})}\\
&\le C_{6}\|Y_{n}\|_{L^{\infty}(\RR\times S^{1};\RR^{m})}+C_{5}\bigl\|(d\widetilde{\scr{F}})_{0}(Y_{n})\bigr\|_{L^{p}(\RR\times S^{1};\RR^{m})}\text{。}
\end{aligned}
\]
第一项由\cref{thm:lift-exp-decay}趋于 $0$。对第二项，我们提醒一下 $\widetilde{\ell}_{n}=\exp_{w_{n}}Y_{n}$，从而 $\widetilde{\scr{F}}(Y_{n})=\scr{F}(\ell_{n})=0$，因为 $\ell_{n}$ 是一条\trFloer\trTrajectory。让我们为 $\widetilde{\scr{F}}$ 写出一个\trTaylor{}型的公式，它类似于我们为了对 $F$ 应用\trNewton--\trPicard{}方法（\cref{thm:abstract-newton-picard}）而写出的那个公式。我们有
\[
\widetilde{\scr{F}}(Y)=\widetilde{\scr{F}}(0)+(d\widetilde{\scr{F}})_{0}(Y)+\scr{N}(Y)\text{，}
\]
其中 $\scr{N}:W^{1,p}(\RR\times S^{1};\RR^{m})\longrightarrow L^{p}(\RR\times S^{1};\RR^{m})$ 满足
\[
\scr{N}(0)=0\qquad\text{并且}\qquad(d\scr{N})_{Z}=(d\widetilde{\scr{F}})_{Z}-(d\widetilde{\scr{F}})_{0}\text{。}
\]
把它应用于 $Y=Y_{n}$，这就给出
\[
0=\widetilde{\scr{F}}(0)+(d\widetilde{\scr{F}})_{0}(Y_{n})+\scr{N}(Y_{n})\text{。}
\]
这正是我们将用来给出上面不等式中第二项之上界的东西。下面是另一条引理。

\begin{lemma}\label{thm:taylor-remainder-estimate}
存在常数 $C_{7}>0$ 与 $r_{1}>0$，使得对每个 $Y\in W^{1,p}(\RR\times S^{1};\RR^{m})$，\footnote{译注：此型别原书印作 $W^{1,p}(\RR\times S^{1};\RR^{2n})$；按上文 $\scr{N}:W^{1,p}(\RR\times S^{1};\RR^{m})\longrightarrow L^{p}(\RR\times S^{1};\RR^{m})$，本引理所论者即此映射，疑应作 $\RR^{m}$。此处径改。——译者注}若 $\|Y\|_{L^{\infty}}\le r_{1}$，那么
\[
\|\scr{N}(Y)\|_{L^{p}}\le C_{7}\|Y\|_{L^{\infty}}\|Y\|_{W^{1,p}}\text{。}
\]
\end{lemma}

让我们暂时承认这条引理。我们于是有
\[
\bigl\|(d\widetilde{\scr{F}})_{0}(Y_{n})\bigr\|_{L^{p}}=\bigl\|\widetilde{\scr{F}}(0)+\scr{N}(Y_{n})\bigr\|_{L^{p}}\le\bigl\|\widetilde{\scr{F}}(0)\bigr\|_{L^{p}}+\|\scr{N}(Y_{n})\|_{L^{p}}\text{。}
\]
由于 $\|Y_{n}\|_{L^{\infty}}$ 趋于 $0$，我们可以对充分大的 $n$ 应用这条引理。这就给出
\[
\bigl\|(d\widetilde{\scr{F}})_{0}(Y_{n})\bigr\|_{L^{p}}\le\bigl\|\widetilde{\scr{F}}(0)\bigr\|_{L^{p}}+C_{7}\|Y_{n}\|_{L^{\infty}}\|Y_{n}\|_{W^{1,p}}\text{。}
\]
把对 $Y_{n}$ 的 $W^{1,p}$ 范数的上界、对 $L(\zeta_{n})$ 的 $L^{p}$ 范数的上界与刚刚得到的那个结合起来，我们发现：对充分大的 $n$，
\[
\begin{aligned}
\|Y_{n}\|_{W^{1,p}}&\le(C_{1}+C_{3})\|\chi_{n}\|_{W^{1,p}}+C_{2}\|L(\zeta_{n})\|_{L^{p}}\\
&\le(C_{1}+C_{3})\|\chi_{n}\|_{W^{1,p}}+C_{8}\|Y_{n}\|_{L^{\infty}}+C_{9}\bigl\|(d\widetilde{\scr{F}})_{0}(Y_{n})\bigr\|_{L^{p}}\\
&\le(C_{1}+C_{3})\|\chi_{n}\|_{W^{1,p}}+C_{8}\|Y_{n}\|_{L^{\infty}}\\
&\qquad\qquad+C_{9}\bigl\|\widetilde{\scr{F}}(0)\bigr\|_{L^{p}}+C_{10}\|Y_{n}\|_{L^{\infty}}\|Y_{n}\|_{W^{1,p}}\text{，}
\end{aligned}
\]
这就蕴含
\[
\|Y_{n}\|_{W^{1,p}}\le\frac{(C_{1}+C_{3})\|\chi_{n}\|_{W^{1,p}}+C_{8}\|Y_{n}\|_{L^{\infty}}+C_{9}\|\widetilde{\scr{F}}(0)\|_{L^{p}}}{1-C_{10}\|Y_{n}\|_{L^{\infty}}}\text{。}
\]
\end{proof}

我们尚须证明\cref{thm:taylor-remainder-estimate}。

\begin{proof}[\cref{thm:taylor-remainder-estimate}的证明]
我们有
\[
\begin{aligned}
\|\scr{N}(Y)\|_{L^{p}}=\|\scr{N}(Y)-\scr{N}(0)\|_{L^{p}}&=\Bigl\|\int_{0}^{1}\frac{\partial}{\partial\sigma}\scr{N}(\sigma Y)\,d\sigma\Bigr\|_{L^{p}}\\
&=\Bigl\|\int_{0}^{1}(d\scr{N})_{\sigma Y}(Y)\,d\sigma\Bigr\|_{L^{p}}\\
&\le\Bigl\|\int_{0}^{1}\bigl\|(d\scr{N})_{\sigma Y}(Y)\bigr\|\,d\sigma\Bigr\|_{L^{p}}\\
&=\Bigl\|\int_{0}^{1}\bigl\|\bigl((d\widetilde{\scr{F}})_{\sigma Y}-(d\widetilde{\scr{F}})_{0}\bigr)Y\bigr\|\,d\sigma\Bigr\|_{L^{p}}\text{。}
\end{aligned}
\]
我们取定一个对子 $(s,t)\in\RR\times S^{1}$。我们在验证\cref{thm:composition-rho-derivative-estimate}的条件 (3) 时（见\cref{sec:proof-second}），曾对\trEuclidean{}范数 $\bigl\|\bigl((d\widetilde{\scr{F}})_{X}-(d\widetilde{\scr{F}})_{0}\bigr)(v)(s,t)\bigr\|$ 得到过一些估计。对某个 $r_{1}>0$，这些估计在 $\|X(s,t)\|\le r_{1}$ 时成立。我们重新启用这些估计，但首先注意到：虽然它们是对 $w_{\rho}=w_{\operatorname{Id},\rho}$ 得到的，它们对 $w_{n}=w_{\nu,\rho_{n}}$ 也成立。于是我们必须证明：当 $\|X(s,t)\|<r_{1}$ 时，存在一个 $K>0$，使得
\[
\begin{aligned}
\bigl\|\bigl((d\widetilde{\scr{F}})_{X}-(d\widetilde{\scr{F}})_{0}\bigr)v(s,t)\bigr\|&\le K\,\|X(s,t)\|\,\|v(s,t)\|\\
&\qquad+\Bigl\|\frac{\partial X}{\partial s}(s,t)\Bigr\|\,\|v(s,t)\|+\Bigl\|\frac{\partial X}{\partial t}(s,t)\Bigr\|\,\|v(s,t)\|\\
&\qquad+\|X(s,t)\|\,\Bigl\|\frac{\partial v}{\partial s}(s,t)\Bigr\|+\|X(s,t)\|\,\Bigl\|\frac{\partial v}{\partial t}(s,t)\Bigr\|\text{。}
\end{aligned}
\]
我们把 $X$ 换成 $\sigma Y$、把 $v$ 换成 $Y$，便发现：对一个不依赖于 $n$ 的常数 $C_{7}$，当 $\|Y\|_{L^{\infty}}<r_{1}$ 时，我们有
\[
\begin{aligned}
\|\scr{N}(Y)\|_{L^{p}}&\le\Bigl\|\int_{0}^{1}C_{7}\sigma\Bigl(\|Y(s,t)\|^{2}\\
&\qquad\qquad+\|Y(s,t)\|\,\Bigl\|\frac{\partial Y}{\partial s}(s,t)\Bigr\|+\|Y(s,t)\|\,\Bigl\|\frac{\partial Y}{\partial t}(s,t)\Bigr\|\Bigr)\,d\sigma\Bigr\|_{L^{p}}\\
&\le C_{7}\Bigl\|\|Y(s,t)\|\Bigl(\|Y(s,t)\|+\Bigl\|\frac{\partial Y}{\partial s}(s,t)\Bigr\|+\Bigl\|\frac{\partial Y}{\partial t}(s,t)\Bigr\|\Bigr)\Bigr\|_{L^{p}}\\
&\le C_{7}\|Y\|_{L^{\infty}}\Bigl\|\Bigl(\|Y\|+\Bigl\|\frac{\partial Y}{\partial s}\Bigr\|+\Bigl\|\frac{\partial Y}{\partial t}\Bigr\|\Bigr)\Bigr\|_{L^{p}}\\
&\le C_{7}\|Y\|_{L^{\infty}}\|Y\|_{W^{1,p}}\text{，}
\end{aligned}
\]
这就证明了\cref{thm:taylor-remainder-estimate}。
\end{proof}
